\documentclass[11pt]{article}

\usepackage[margin=1in]{geometry}
\usepackage{amsmath,amssymb,amsthm,mathtools}
\usepackage{enumitem}

\newtheorem{theorem}{Theorem}
\newtheorem{lemma}[theorem]{Lemma}
\newtheorem{proposition}[theorem]{Proposition}
\newtheorem{remark}[theorem]{Remark}

\DeclareMathOperator{\Sh}{Sh}
\DeclareMathOperator{\Lan}{Lan}
\DeclareMathOperator{\colim}{colim}
\DeclareMathOperator{\Face}{Face}
\DeclareMathOperator{\Prin}{Prin}
\DeclareMathOperator{\Spec}{Spec}
\DeclareMathOperator{\R}{R}

\begin{document}

\title{Cohomological full faithfulness for the birhizomic affine fan}
\author{}
\date{}
\maketitle

\section*{The result}

We prove that the answer is yes.  Throughout, sheaves on \(A_P\) and \(B_P\) mean sheaves on the sites
\(A_P^{\mathrm{op}}\) and \(B_P^{\mathrm{op}}\).  Since the topology on \(B_P^{\mathrm{op}}\) is trivial, a sheaf on \(B_P\) is the same thing as a presheaf on \(B_P^{\mathrm{op}}\), equivalently a covariant functor
\[
B_P \longrightarrow \mathbf{Ab}.
\]

\begin{theorem}
Let \(P\) be an integral, saturated, sharp commutative monoid.  Let
\[
\pi^{-1} : \Sh(B_P) \longrightarrow \Sh(A_P)
\]
be the functor obtained from the inclusion \(B_P^{\mathrm{op}}\subset A_P^{\mathrm{op}}\) by left Kan extension and sheafification, and let \(\pi_*\) be its right adjoint.  Then, for every presheaf of abelian groups \(F\) on \(B_P^{\mathrm{op}}\), the adjunction morphism
\[
F \longrightarrow \R\pi_*\pi^{-1}F
\]
is a quasi-isomorphism.
\end{theorem}

The proof uses the standard affine comparison theorem for valuative, equivalently birhizomic, fans.  This theorem is the affine form of Kato's cohomological contractibility theorem for valuative fans; see Kato--Nakayama--Usui \cite[Section 3]{KNU}, Kato \cite[Sections 1--3]{Kato}, and Ogus \cite[Chapter V, Sections 1--2]{Ogus}.  In the language of sites it is also an application of the locally-aspherical comparison theorem of SGA \(4\), Exposé IX; see \cite[Exposé IX, Theorem 5.7]{SGA4}.

\section{Faces and the left Kan extension}

A face of \(P\) is a submonoid \(S\subset P\) such that, whenever \(a+b\in S\), then \(a,b\in S\).  If \(x\in P\), define the principal face generated by \(x\) by
\[
S_x := \{p\in P \mid p\leq n x \text{ for some } n\geq 0\}.
\]
Equivalently, \(S_x=\{p\in P\mid p\prec x\}\).

For a face \(S\subset P\), write \(P/S\) for the sharp monoid obtained from \(P\) by making every element of \(S\) invertible and then sharpifying.  Thus \(P\to P/S\) is initial among homomorphisms from \(P\) which send every element of \(S\) to \(0\).

\begin{lemma}
For every \(x\in P\), the monoid \(P[-x]^\sharp\) is canonically isomorphic to \(P/S_x\).
\end{lemma}

\begin{proof}
Let \(N\) be a sharp monoid and let \(\varphi:P\to N\) be a homomorphism with \(\varphi(x)=0\).  If \(p\in S_x\), then \(p+r=n x\) for some \(r\in P\) and some \(n\geq 0\).  Hence
\[
\varphi(p)+\varphi(r)=n\varphi(x)=0.
\]
Since \(N\) is sharp, \(\varphi(p)=0\).  Thus every map killing \(x\) kills all of \(S_x\).  Conversely \(x\in S_x\), so every map killing \(S_x\) kills \(x\).  The two universal properties are therefore the same.
\end{proof}

Hence, up to equivalence, the category \(B_P\) may be identified with the full subcategory of \(A_P\) consisting of the quotients \(P/S\) for \(S\) a principal face.  There is a unique morphism
\[
P/S \longrightarrow P/T
\]
over \(P\) if and only if \(S\subset T\).

For \(Q\in A_P\), let
\[
K_Q := \ker(P\to Q)=\{p\in P\mid p\mapsto 0\text{ in }Q\}.
\]
This is a face of \(P\).

Let \(F:B_P\to \mathbf{Ab}\) be a presheaf on \(B_P^{\mathrm{op}}\).  Define a functor
\[
\overline F:\Face(P)\longrightarrow \mathbf{Ab}
\]
by
\[
\overline F(S):=\colim_{S_x\subset S} F(P/S_x),
\]
where the colimit is over the filtered category of principal faces contained in \(S\).  If \(S=S_x\) is principal, then \(S_x\) is final in that indexing category, so
\[
\overline F(S_x)\cong F(P/S_x).
\]

\begin{lemma}
Let \(i:B_P\hookrightarrow A_P\) denote the inclusion.  The presheaf left Kan extension of \(F\) along \(i\) is the presheaf
\[
Q\longmapsto \overline F(K_Q).
\]
Consequently
\[
\pi^{-1}F \cong a\bigl(Q\mapsto \overline F(K_Q)\bigr),
\]
where \(a\) denotes sheafification for the birhizomic topology.
\end{lemma}

\begin{proof}
By the formula for a left Kan extension,
\[
(\Lan_iF)(Q)=\colim_{(P/S_x\to Q)}F(P/S_x).
\]
A morphism \(P/S_x\to Q\) over \(P\) exists if and only if \(S_x\subset K_Q\), and in that case it is unique.  Therefore the indexing category is precisely the category of principal faces contained in \(K_Q\), giving
\[
(\Lan_iF)(Q)=\colim_{S_x\subset K_Q}F(P/S_x)=\overline F(K_Q).
\]
Sheafifying gives the stated formula for \(\pi^{-1}F\).
\end{proof}

\section{The affine birhizomic comparison theorem}

We use the following standard affine comparison theorem.

\begin{theorem}[Affine valuative/birhizomic comparison]\label{thm:affine-comparison}
Let \(M\) be an integral, saturated, sharp monoid, and let
\[
E:\Face(M)\longrightarrow \mathbf{Ab}
\]
be a covariant functor.  Let \(E^{\mathrm{br}}\) be the sheaf on \(A_M^{\mathrm{op}}\) associated to the presheaf
\[
R\longmapsto E\bigl(\ker(M\to R)\bigr).
\]
Then the canonical map
\[
E(\{0\}) \longrightarrow \R\Gamma(A_M^{\mathrm{op}},E^{\mathrm{br}})
\]
is an isomorphism in \(D(\mathbf{Ab})\).

More generally, if \(S\subset M\) is a face, then applying the same statement to \(M/S\) gives
\[
E(S) \simeq \R\Gamma(A_{M/S}^{\mathrm{op}},E_S^{\mathrm{br}}),
\]
where \(E_S\) is the induced functor on faces of \(M/S\), i.e. on faces of \(M\) containing \(S\).
\end{theorem}

\begin{proof}[Reference and translation]
This is the affine form of the cohomological contractibility theorem for the valuative fan of an affine fan.  In Kato's language, \(\Spec(M)_{\mathrm{val}}\to \Spec(M)\) is locally aspherical over every stratum.  The birhizomic topology defined in the question is precisely the site-theoretic formulation of the valuative topology on affine fans: birhizomic morphisms are the finite presentation morphisms which satisfy the valuative lifting criterion, and birhizomic covers are the covers by valuative subdivisions.  See Kato \cite[Sections 1--3]{Kato}, Kato--Nakayama--Usui \cite[Section 3]{KNU}, and Ogus \cite[Chapter V, Sections 1--2]{Ogus}.  The passage from local asphericity to the displayed derived global-section statement is exactly the comparison theorem for locally aspherical morphisms of sites, SGA \(4\), Exposé IX, Theorem 5.7 \cite{SGA4}.
\end{proof}

\section{Proof of the theorem}

Let \(F:B_P\to \mathbf{Ab}\) be a presheaf on \(B_P^{\mathrm{op}}\).  We must show that
\[
F \longrightarrow \R\pi_*\pi^{-1}F
\]
is a quasi-isomorphism.

Because the topology on \(B_P^{\mathrm{op}}\) is trivial, it is enough to check this objectwise.  Thus fix \(x\in P\), and put
\[
S:=S_x,\qquad b_x:=P/S=P[-x]^\sharp.
\]

The right adjoint \(\pi_*\) is restriction to \(B_P\).  Therefore, for a sheaf \(G\) on \(A_P^{\mathrm{op}}\),
\[
(\R\pi_*G)(b_x)
\simeq
\R\Gamma\bigl((A_P^{\mathrm{op}})/b_x,G\bigr).
\]
This is the standard formula identifying the derived functors of evaluation at an object with cohomology of the slice site; see SGA \(4\), Exposé V, Proposition 5.1 \cite{SGA4}.

Now
\[
(A_P^{\mathrm{op}})/b_x \simeq A_{b_x}^{\mathrm{op}}=A_{P/S}^{\mathrm{op}},
\]
because an object of the slice is a morphism \(b_x\to Q\) in \(A_P\), equivalently a monoid homomorphism \(P/S\to Q\).

By the computation of the left Kan extension above,
\[
\pi^{-1}F
\cong
a\bigl(Q\mapsto \overline F(K_Q)\bigr).
\]
Restricting to the slice \(A_{P/S}^{\mathrm{op}}\), this becomes the sheaf associated to the presheaf
\[
Q\longmapsto \overline F(K_Q),
\]
where now \(K_Q\) is viewed as a face of \(P\) containing \(S\).  Equivalently, it is the sheaf \(E^{\mathrm{br}}\) associated to the functor
\[
E:\Face(P/S)\longrightarrow \mathbf{Ab},
\qquad
T/S\longmapsto \overline F(T).
\]

Applying Theorem \ref{thm:affine-comparison} to the monoid \(P/S\), we obtain
\[
\R\Gamma\bigl(A_{P/S}^{\mathrm{op}},(\pi^{-1}F)|_{A_{P/S}}\bigr)
\simeq
E(\{0\})
=
\overline F(S).
\]
Since \(S=S_x\) is principal, \(S\) is final among the principal faces contained in \(S\), and hence
\[
\overline F(S)\cong F(P/S)=F(b_x).
\]
Combining these identifications gives
\[
(\R\pi_*\pi^{-1}F)(b_x)\simeq F(b_x),
\]
concentrated in degree \(0\).

The construction is functorial in \(b_x\), and the degree \(0\) map is exactly the adjunction unit
\[
F(b_x)\longrightarrow (\pi_*\pi^{-1}F)(b_x).
\]
Thus the adjunction morphism
\[
F\longrightarrow \R\pi_*\pi^{-1}F
\]
is an objectwise quasi-isomorphism.  Since the topology on \(B_P^{\mathrm{op}}\) is trivial, objectwise quasi-isomorphisms are precisely quasi-isomorphisms of complexes of sheaves on \(B_P^{\mathrm{op}}\).  This proves the theorem.

\begin{thebibliography}{9}

\bibitem{SGA4}
M. Artin, A. Grothendieck, and J.-L. Verdier,
\emph{Théorie des topos et cohomologie étale des schémas},
Lecture Notes in Mathematics, Vols. 269, 270, 305, Springer-Verlag, 1972--1973.

\bibitem{Kato}
K. Kato,
\emph{Toric singularities},
Amer. J. Math. \textbf{116} (1994), no. 5, 1073--1099.

\bibitem{KNU}
K. Kato, C. Nakayama, and S. Usui,
\emph{Classifying spaces of degenerating polarized Hodge structures},
Annals of Mathematics Studies, vol. 169, Princeton University Press, 2009.

\bibitem{Ogus}
A. Ogus,
\emph{Lectures on Logarithmic Algebraic Geometry},
Cambridge Studies in Advanced Mathematics, vol. 178, Cambridge University Press, 2018.

\end{thebibliography}

\end{document}
```